\documentclass[11pt]{article}
\usepackage[letterpaper,margin=1in]{geometry}
\usepackage[T1]{fontenc}
\usepackage{lmodern}
\usepackage{amsmath,amssymb,amsthm}
\usepackage[colorlinks=true,linkcolor=blue,urlcolor=blue,citecolor=blue]{hyperref}
\usepackage{microtype}

\newtheorem*{theorem}{Theorem}
\newcommand{\R}{\mathbb R}

\title{The Reciprocal-Sum Neumann Inequality:\\
Literature Resolution of an Open Problem in arXiv:1604.05072}
\author{Run \texttt{fa\_banach\_001} --- literature-status note}
\date{August 11, 2026}

\begin{document}
\maketitle

\noindent\textbf{Status.} \texttt{literature\_already\_answered}. This note
identifies an exact later solution. It does not claim a new theorem.

\section*{The source problem}
Brasco and De Philippis use the convention
\[
 0=\mu_1(\Omega)<\mu_2(\Omega)\leq\mu_3(\Omega)\leq\cdots
\]
for the Neumann spectrum.  In the subsection ``A quick overview of the
Neumann spectrum,'' they ask whether the higher-dimensional analogue of
Szeg\H{o}'s planar inequality is valid:
\begin{equation}
 \frac{1}{|\Omega|^{2/N}}
 \sum_{k=2}^{N+1}\frac{1}{\mu_k(\Omega)}
 \geq
 \frac{1}{|B|^{2/N}}
 \sum_{k=2}^{N+1}\frac{1}{\mu_k(B)}.                 \tag{S}
\end{equation}
Here $B\subset\R^N$ is a ball.  Because its first nonzero eigenvalue has
multiplicity $N$, the right side is
$N/(|B|^{2/N}\mu_2(B))$.  The survey says explicitly that the validity of
(S) was open \cite{BD}.

\section*{Exact later theorem}
He, Li, and Tang number the Neumann eigenvalues instead as
\[
 0=\mu_0(\Omega)<\mu_1(\Omega)\leq\mu_2(\Omega)\leq\cdots.
\]
Their main result is the following \cite[Theorem 1.2]{HLT}.

\begin{theorem}
Let $N\geq2$, let $\Omega\subset\R^N$ be a bounded domain with smooth
boundary, and let $B$ be a ball with $|B|=|\Omega|$. Then
\[
 \sum_{j=1}^{N}\frac{1}{\mu_j(\Omega)}
 \geq \frac{N}{\mu_1(B)}.
\]
Equality holds if and only if $\Omega$ is a ball.
\end{theorem}

This is the Ashbaugh--Benguria conjecture.  Li and Wang subsequently record
that the proof remains valid for every bounded Lipschitz domain and state the
result in that form \cite[Theorem 1.1]{LW}.  They also prove stronger
quantitative stability estimates for the same deficit.

\section*{Translation to the source notation}
The index shift is the only spectral difference:
the later $\mu_j$ is the source's $\mu_{j+1}$.  To check the normalization,
let $\mathbb B$ be the unit ball, let
\[
 R_\Omega=\left(\frac{|\Omega|}{|\mathbb B|}\right)^{1/N},
 \qquad \widetilde\mu_j(\Omega)=R_\Omega^2\mu_j(\Omega).
\]
Li and Wang state the theorem as
\[
 \sum_{j=1}^{N}\frac{1}{\widetilde\mu_j(\Omega)}
 \geq \frac{N}{\mu_1(\mathbb B)}.                  \tag{1}
\]
Since $R_\Omega^{-2}=(|\mathbb B|/|\Omega|)^{2/N}$,
division of (1) by $|\mathbb B|^{2/N}$ gives
\[
 \frac{1}{|\Omega|^{2/N}}
 \sum_{j=1}^{N}\frac{1}{\mu_j(\Omega)}
 \geq
 \frac{N}{|\mathbb B|^{2/N}\mu_1(\mathbb B)}.
\]
Scale invariance and the $N$-fold equality
$\mu_1(B)=\cdots=\mu_N(B)$ turn the right side into
\[
 \frac{1}{|B|^{2/N}}\sum_{j=1}^{N}\frac{1}{\mu_j(B)}.
\]
After shifting $j$ to $k=j+1$, this is exactly (S).  The equality statement
also transfers: equality occurs precisely for balls.

For a disconnected Lipschitz open set, the source convention supplies
additional zero eigenvalues; the reciprocal expression is then infinite (or
is excluded from the finite formulation).  Thus the substantive connected
case is exactly the domain case covered above.

\section*{Proof provenance and scope}
The later proof is not a citation-level restatement.  It couples the $N$
trial functions transplanted from the first eigenspace of the ball, applies a
trace-form Ritz/Hersch variational principle, and proves the needed matrix
comparison by exploiting a trace-free angular-imbalance matrix.  Equality in
the matrix and raywise rearrangement steps forces the domain to be a ball
\cite{HLT}.

This packet resolves only the displayed reciprocal-sum Neumann problem.  The
same source also asks for the sharp topology-free upper bound on
$P(\Omega)\sigma_2(\Omega)$ for planar Steklov eigenvalues.  That is a
separate question and is not claimed to follow from the Neumann theorem.

\begin{thebibliography}{9}
\bibitem{BD}
L. Brasco and G. De Philippis,
\emph{Spectral inequalities in quantitative form},
arXiv:1604.05072 (2016), Section 4.

\bibitem{HLT}
Y. He, Y. Li, and Q. Tang,
\emph{A proof of the Ashbaugh--Benguria conjecture for reciprocal sums of
Neumann eigenvalues}, arXiv:2606.08271 (2026).

\bibitem{LW}
Y. Li and Z. Wang,
\emph{Spectral and Geometric Stability for the Reciprocal Sum of Neumann
Eigenvalues}, arXiv:2607.19008 (2026).
\end{thebibliography}

\end{document}
